\section{Kết quả Thực nghiệm và Thảo luận Khoa học}

\subsection{Tổng quan Kết quả Thực nghiệm qua Hai Kịch bản}
Bảng~\ref{tab:main_benchmark_all} và Bảng~\ref{tab:main_benchmark_filtered} trình bày kết quả đánh giá chi tiết của toàn bộ các cấu hình thuật toán trên tập kiểm thử chuẩn quốc tế \textsc{resume-job-description-fit} ($1.759$ cặp nhãn trên $71$ truy vấn JD) lần lượt dưới hai kịch bản Setting A (toàn bộ truy vấn) và Setting B (các truy vấn đã lọc theo chuẩn TREC).

% Bảng kết quả tổng thể trên toàn bộ 71 truy vấn (Setting A: Chưa lọc)
\begin{table*}[t]
\centering
\caption{Hiệu năng truy xuất và xếp hạng trên toàn bộ 71 truy vấn JD (Kịch bản Setting A: Đánh giá toàn diện, các truy vấn không có ứng viên phù hợp $\text{rel}\ge 1$ nhận điểm $0.0$) trên tập dữ liệu \texttt{cnamuangtoun/resume-job-description-fit}. $T_{\text{rank}}$ là độ trễ tính điểm và xếp hạng trực tuyến trên CPU cho tập ứng viên $N=25$.}
\label{tab:main_benchmark_all}
\resizebox{\textwidth}{!}{%
\begin{tabular}{lrrrrrrr}
\hline
\textbf{Phương pháp / Kiến trúc} & \textbf{nDCG@5} & \textbf{nDCG@10} & \textbf{MRR} & \textbf{MAP} & \textbf{P@5} & \textbf{Spearman} & \textbf{$T_{\text{rank}}$ ($N=25$)} \\
\hline
Xếp hạng Ngẫu nhiên (Floor) & 0.4532 & 0.4963 & 0.5721 & 0.4941 & 0.4173 & -0.026 & -- \\
BM25 / TF-IDF (Từ vựng) & 0.5532 & 0.5773 & 0.6393 & 0.5930 & 0.4592 & 0.209 & 0.16\,ms \\
Đơn Vector Phẳng MiniLM (Flat) & 0.5889 & 0.6146 & 0.7055 & 0.6210 & 0.4845 & 0.256 & $<0.10$\,ms \\
Chunk-Mean (MiniLM, cửa sổ 128 từ) & 0.6123 & 0.6193 & 0.6982 & 0.6356 & 0.5099 & 0.306 & $<0.10$\,ms \\
Chunk-MaxSim (Tương tác trễ ColBERT) & 0.5758 & 0.5992 & 0.6785 & 0.6119 & 0.4817 & 0.267 & 0.39\,ms \\
Đa Vector Nhận biết Mục (Ours) & 0.6053 & 0.6170 & 0.7091 & 0.6328 & 0.4986 & 0.283 & 0.12\,ms \\
BGE-M3 Phẳng (Ngữ cảnh 8.192 tokens) & 0.5688 & 0.5937 & 0.6555 & 0.6167 & 0.4789 & 0.248 & $<0.10$\,ms \\
BGE-M3 Nhận biết Phân mục & 0.5934 & 0.6086 & 0.7126 & 0.6169 & 0.4901 & 0.257 & 0.13\,ms \\
Dense + BM25 ($L_0$: Không Taxonomy) & 0.6009 & 0.6123 & 0.6801 & 0.6337 & 0.4873 & 0.276 & 0.28\,ms \\
Mô hình Lai Hiệu chuẩn (Calibrated Hybrid) & 0.5998 & 0.6162 & 0.6764 & 0.6366 & 0.4817 & 0.295 & 2.69\,ms \\
RRF + Taxonomy ($k=60$) & 0.5980 & 0.6107 & 0.6870 & 0.6299 & 0.4873 & 0.293 & 2.83\,ms \\
Tái xếp hạng 2 Chặng: Hybrid $\rightarrow$ CE (Top-10) & 0.5513 & 0.5766 & 0.6408 & 0.5767 & 0.4761 & 0.212 & 426.9\,ms \\
Cross-Encoder Toàn phần (ms-marco-MiniLM) & 0.5292 & 0.5531 & 0.6420 & 0.5650 & 0.4592 & 0.167 & 1178.0\,ms \\
Xếp hạng Hoàn hảo Oracle (Ceiling) & 0.8592 & 0.8592 & 0.8592 & 0.8592 & 0.6958 & 0.791 & -- \\
\hline
\end{tabular}%
}
\end{table*}


\subsection{Đặc trưng Cấu trúc Phân mục và Độ bao phủ Ký tự Thực tế}
Bảng~\ref{tab:section_skills_stats} cung cấp cái nhìn định lượng về mức độ phong phú phân mục trên tập dữ liệu gồm $1.051$ CV và $351$ JD độc nhất sau khi xử lý bởi bộ tách đã đóng băng (commit \texttt{41e1d27}).

% Thống kê cấu trúc phân mục thống nhất trên tập dữ liệu benchmark công khai
\begin{table}[t]
\centering
\caption{Thống kê phân bố phân mục trên CV và JD được bóc tách bởi bộ tách phân mục đã đóng băng (commit \texttt{41e1d27}) trên tập dữ liệu chuẩn \texttt{cnamuangtoun/resume-job-description-fit}.}
\label{tab:section_skills_stats}
\resizebox{\textwidth}{!}{%
\begin{tabular}{lrrrr}
\hline
\textbf{Thực thể Phân mục (Section Entity)} & \textbf{Tỷ lệ CV (\%)} & \textbf{Số từ TB / CV} & \textbf{Tỷ lệ JD (\%)} & \textbf{Số từ TB / JD} \\
\hline
Kỹ năng Kỹ thuật / Yêu cầu công việc & 86.4\% & 49.8 & 19.7\% & 250.6 \\
Kinh nghiệm Làm việc / Trách nhiệm chính & 100.0\% & 527.0 & 31.0\% & 205.6 \\
Học vấn \& Bằng cấp & 94.3\% & 38.9 & 8.5\% & 253.8 \\
Tóm tắt Bản thân / Tổng quan hồ sơ & 83.9\% & 86.3 & -- & -- \\
Chế độ Đãi ngộ \& Quyền lợi & -- & -- & 15.5\% & 172.6 \\
Dự án / Portfolio & -- & -- & -- & -- \\
\hline
Đa mục ($\ge 2$ phân mục nhận diện) & 99.6\% & -- & 28.2\% & -- \\
Đa mục ($\ge 3$ phân mục nhận diện) & 94.3\% & -- & 7.0\% & -- \\
Bao phủ Kỹ năng Taxonomy & 96.4\% & 10.3 kỹ năng/CV & 87.3\% & 4.4 kỹ năng/JD \\
\hline
\end{tabular}%
}
\end{table}


Kết quả cho thấy sự phân hóa cấu trúc sâu sắc giữa hai loại tài liệu:
\begin{itemize}
    \item \textbf{Cấu trúc Hồ sơ Ứng viên (CV)}: $86.4\%$ CV chứa mục Kỹ năng Kỹ thuật riêng, và $100.0\%$ CV có mục Kinh nghiệm Làm việc. Đặc biệt, $99.6\%$ CV sở hữu từ hai phân mục trở lên, và $94.3\%$ sở hữu từ ba phân mục. Đo đạc độ bao phủ ký tự chức năng khẳng định \textbf{$95.32\%$ số ký tự của CV rơi vào các phân mục chức năng cốt lõi} (kỹ năng, kinh nghiệm, học vấn), trong khi phần mở đầu (preamble) chỉ chiếm trung vị $2.00\%$. Điều này chứng minh CV là loại tài liệu có cấu trúc cao.
    \item \textbf{Cấu trúc Mô tả Công việc (JD)}: Ngược lại với CV, các bài đăng tuyển dụng thể hiện tính phi cấu trúc rõ rệt: chỉ có $19.7\%$ JD có mục kỹ năng riêng và $31.0\%$ có mục trách nhiệm công việc riêng. Phần lời tự sự mở đầu (giới thiệu công ty, tầm nhìn) chiếm tới trung vị \textbf{$67.72\%$ tổng số ký tự của JD}, và chỉ có $30.71\%$ ký tự rơi vào các phân mục chức năng. Sự bất đối xứng cấu trúc này lý giải vì sao cơ chế ghép cặp chéo (Cross-Section Pairing) lại mang tính sống còn đối với hệ thống.
\end{itemize}

% Bảng kết quả trên các truy vấn chuẩn TREC đã lọc (Setting B: N=61 truy vấn có ít nhất 1 ứng viên phù hợp rel >= 1)
\begin{table*}[t]
\centering
\caption{Hiệu năng truy xuất trên tập truy vấn đã lọc theo chuẩn TREC có ít nhất một ứng viên phù hợp (Kịch bản Setting B: $N=61$ truy vấn tại ngưỡng $\text{rel} \ge 1$). $T_{\text{rank}}$ là độ trễ tính điểm trực tuyến trên CPU cho tập ứng viên $N=25$.}
\label{tab:main_benchmark_filtered}
\resizebox{\textwidth}{!}{%
\begin{tabular}{lrrrrrrr}
\hline
\textbf{Phương pháp / Kiến trúc} & \textbf{nDCG@5} & \textbf{nDCG@10} & \textbf{MRR} & \textbf{MAP} & \textbf{P@5} & \textbf{Spearman} & \textbf{$T_{\text{rank}}$ ($N=25$)} \\
\hline
Xếp hạng Ngẫu nhiên (Floor) & 0.5280 & 0.5743 & 0.6639 & 0.5741 & 0.4877 & -0.026 & -- \\
BM25 / TF-IDF (Từ vựng) & 0.6439 & 0.6720 & 0.7441 & 0.6903 & 0.5344 & 0.209 & 0.16\,ms \\
Đơn Vector Phẳng MiniLM (Flat) & 0.6854 & 0.7153 & 0.8211 & 0.7228 & 0.5639 & 0.256 & $<0.10$\,ms \\
Chunk-Mean (MiniLM, cửa sổ 128 từ) & 0.7127 & 0.7208 & 0.8126 & 0.7398 & 0.5934 & 0.306 & $<0.10$\,ms \\
Chunk-MaxSim (Tương tác trễ ColBERT) & 0.6702 & 0.6974 & 0.7898 & 0.7122 & 0.5607 & 0.267 & 0.39\,ms \\
Đa Vector Nhận biết Mục (Ours) & 0.7046 & 0.7182 & 0.8253 & 0.7365 & 0.5803 & 0.283 & 0.12\,ms \\
BGE-M3 Phẳng (Ngữ cảnh 8.192 tokens) & 0.6620 & 0.6910 & 0.7630 & 0.7178 & 0.5574 & 0.248 & $<0.10$\,ms \\
BGE-M3 Nhận biết Phân mục & 0.6907 & 0.7084 & 0.8294 & 0.7181 & 0.5705 & 0.257 & 0.13\,ms \\
Dense + BM25 ($L_0$: Không Taxonomy) & 0.6994 & 0.7127 & 0.7916 & 0.7376 & 0.5672 & 0.276 & 0.28\,ms \\
Mô hình Lai Hiệu chuẩn (Calibrated Hybrid) & 0.6981 & 0.7173 & 0.7873 & 0.7410 & 0.5607 & 0.295 & 2.69\,ms \\
RRF + Taxonomy ($k=60$) & 0.6960 & 0.7109 & 0.7996 & 0.7332 & 0.5672 & 0.293 & 2.83\,ms \\
Tái xếp hạng 2 Chặng: Hybrid $\rightarrow$ CE (Top-10) & 0.6416 & 0.6711 & 0.7459 & 0.6712 & 0.5541 & 0.212 & 426.9\,ms \\
Cross-Encoder Toàn phần (ms-marco-MiniLM) & 0.6160 & 0.6437 & 0.7473 & 0.6577 & 0.5344 & 0.167 & 1178.0\,ms \\
Xếp hạng Hoàn hảo Oracle (Ceiling) & 1.0000 & 1.0000 & 1.0000 & 1.0000 & 0.8098 & 0.791 & -- \\
\hline
\end{tabular}%
}
\end{table*}


\subsection{Kiểm định Ý nghĩa Thống kê trên Họ Giả thuyết Chính}
Bảng~\ref{tab:statistical_tests} báo cáo toàn bộ kết quả kiểm định thống kê ghép cặp trên $15$ giả thuyết cốt lõi thuộc các câu hỏi nghiên cứu RQ1--RQ4 ($N=71$ truy vấn, Setting A), bao gồm hiệu số trung bình $\bar{\Delta}$, khoảng tin cậy Bootstrap 95\% CI tính trên $10.000$ mẫu lặp, giá trị $p$ thô của kiểm định $t$, và giá trị $p$ sau hiệu chỉnh đa so sánh Holm--Bonferroni.

% Bảng kiểm định thống kê ghép cặp Bootstrap 95% CI và Holm-Bonferroni
\begin{table*}[t]
\centering
\caption{Kiểm định Thống kê Ý nghĩa Ghép cặp Đầy đủ trên Họ Giả thuyết Chính ($N=71$ truy vấn, Setting A). Khoảng tin cậy Bootstrap 95\% CI tính trên $10{,}000$ lần lấy mẫu lặp ghép cặp. Đa kiểm định được kiểm soát qua thủ tục Holm-Bonferroni ($\alpha=0.05$).}
\label{tab:statistical_tests}
\resizebox{\textwidth}{!}{%
\begin{tabular}{lllrrrrrc}
\hline
\textbf{RQ} & \textbf{Cặp so sánh ($A$ vs. $B$)} & \textbf{Thước đo} & \textbf{Hiệu TB $\bar{\Delta}$} & \textbf{Bootstrap 95\% CI} & \textbf{Thống kê $t$} & \textbf{$p$ thô} & \textbf{$p$ hiệu chỉnh} & \textbf{Ý nghĩa ($\alpha=0.05$)} \\
\hline
RQ1 & Section-Aware vs. Single-Vector & nDCG@10 & +0.0024 & [-0.0296, +0.0333] & 0.15 & 0.8811 & 1.0000 & Không \\
RQ1 & Section-Aware vs. Single-Vector & MAP & +0.0118 & [-0.0124, +0.0393] & 0.88 & 0.3819 & 1.0000 & Không \\
RQ1 & BGE-M3 Section vs. BGE-M3 Flat & nDCG@10 & +0.0149 & [-0.0202, +0.0496] & 0.83 & 0.4107 & 1.0000 & Không \\
RQ1 & BGE-M3 Section vs. BGE-M3 Flat & MAP & +0.0002 & [-0.0319, +0.0325] & 0.01 & 0.9891 & 1.0000 & Không \\
RQ1 & Chunk-Mean vs. Single-Vector & nDCG@10 & +0.0047 & [-0.0310, +0.0397] & 0.26 & 0.7983 & 1.0000 & Không \\
RQ1 & Chunk-MaxSim vs. Single-Vector & nDCG@10 & -0.0154 & [-0.0419, +0.0095] & -1.15 & 0.2524 & 1.0000 & Không \\
\hline
RQ2 & Single-Vector vs. BM25 & nDCG@10 & +0.0373 & [+0.0018, +0.0736] & 2.01 & 0.0480 & 0.5762 & Không \\
RQ2 & Single-Vector vs. BM25 & MAP & +0.0280 & [-0.0081, +0.0656] & 1.47 & 0.1464 & 1.0000 & Không \\
RQ2 & Dense+BM25 vs. Single-Vector & MAP & +0.0127 & [-0.0181, +0.0476] & 0.75 & 0.4566 & 1.0000 & Không \\
RQ2 & RRF+Taxonomy vs. Dense+BM25 & MAP & -0.0039 & [-0.0233, +0.0119] & -0.42 & 0.6728 & 1.0000 & Không \\
\hline
RQ3 & Calibrated Hybrid vs. Dense+BM25 & MAP & +0.0029 & [-0.0055, +0.0127] & 0.61 & 0.5430 & 1.0000 & Không \\
RQ3 & Calibrated Hybrid vs. Dense+BM25 & nDCG@10 & +0.0039 & [-0.0070, +0.0157] & 0.67 & 0.5051 & 1.0000 & Không \\
\hline
RQ4 & Calibrated Hybrid vs. Cross-Encoder & nDCG@10 & +0.0632 & [+0.0218, +0.1049] & 2.97 & 0.0041 & 0.0586 & Không \\
RQ4 & Calibrated Hybrid vs. Cross-Encoder & MAP & +0.0716 & [+0.0270, +0.1196] & 2.98 & 0.0039 & 0.0586 & Không \\
RQ4 & Calibrated Hybrid vs. 2-Stage Top10 & nDCG@10 & +0.0396 & [+0.0061, +0.0770] & 2.16 & 0.0342 & 0.4440 & Không \\
\hline
\end{tabular}%
}
\end{table*}


\subsection{Trả lời RQ1: Độ mịn Đa Vector Nhận biết Mục so với Đơn Vector Phẳng}
\label{sec:rq1}
RQ1 kiểm tra liệu việc bóc tách phân mục chức năng có mang lại lợi thế xếp hạng vượt trội so với việc nén toàn bộ tài liệu vào một vector đơn lẻ hay không. Bảng~\ref{tab:section_variants} trình bày phân tích phân rã (ablation) giữa các biến thể nhận biết phân mục và các phương pháp chia lát (chunking).

% Bảng phân tích các biến thể Section-Aware và Chunking
\begin{table}[t]
\centering
\caption{So sánh các biến thể Nhận biết Phân mục với Đơn Vector và Chunking, kèm báo cáo tỷ lệ Fallback ($T_{\text{rank}}$ đo trên tập ứng viên $N=25$).}
\label{tab:section_variants}
\resizebox{\columnwidth}{!}{%
\begin{tabular}{lrrrrr}
\hline
\textbf{Biến thể / Cấu hình} & \textbf{nDCG@10} & \textbf{MRR} & \textbf{MAP} & \textbf{Fallback \%} & \textbf{$T_{\text{rank}}$ ($N=25$)} \\
\hline
Đơn Vector Phẳng (256 tokens) & 0.6146 & 0.7055 & 0.6210 & 0.0\% & $<0.10$\,ms \\
Chunk-Mean (128w, stride 64) & 0.6193 & 0.6982 & 0.6356 & 0.0\% & $<0.10$\,ms \\
Chunk-MaxSim (Kiểu ColBERT) & 0.5992 & 0.6785 & 0.6119 & 0.0\% & 0.39\,ms \\
Section-Aware (Biến thể 1: Thuần 100\%) & 0.6170 & 0.7091 & 0.6328 & 0.00\% & 0.12\,ms \\
Section-Aware (Biến thể 2: Có Fallback) & 0.6170 & 0.7091 & 0.6328 & 0.00\% & 0.12\,ms \\
Section-Aware (Biến thể 3: Trộn 0.65/0.35) & 0.6172 & 0.7038 & 0.6321 & 0.00\% & 0.12\,ms \\
BGE-M3 Phẳng (Ngữ cảnh 8.192 tokens) & 0.5937 & 0.6555 & 0.6167 & 0.0\% & $<0.10$\,ms \\
BGE-M3 Nhận biết Phân mục & 0.6086 & 0.7126 & 0.6169 & 0.00\% & 0.13\,ms \\
\hline
\end{tabular}%
}
\end{table}


\noindent\textbf{1. Hiệu năng Xếp hạng và Độ tin cậy của Bộ tách}:
Trên tập kiểm thử đầy đủ (Bảng~\ref{tab:main_benchmark_all}), mô hình Đa Vector Nhận biết Mục (Section-Aware Dense) đạt nDCG@10 là $0.6170$, MRR là $0.7091$ và MAP là $0.6328$, so với $0.6146$ nDCG@10, $0.7055$ MRR và $0.6210$ MAP của Đơn Vector Phẳng (Single-Vector Dense). Trên các truy vấn đã lọc (Setting B, Bảng~\ref{tab:main_benchmark_filtered}), Section-Aware đạt nDCG@10 là $0.7182$ và MAP là $0.7365$ (so với $0.7153$ và $0.7228$ của Single-Vector). Tỷ lệ rơi vào cơ chế dự phòng (fallback rate) của bộ tách đã đóng băng trên toàn bộ $1.759$ cặp kiểm thử đạt chính xác \textbf{$0.00\%$}, khẳng định lỗi dính chữ tiêu đề đã được khắc phục hoàn toàn. Do đó, Biến thể 1 (Thuần phân mục) và Biến thể 2 (Có Fallback) đạt kết quả trùng khớp nhau.

\noindent\textbf{2. Kết luận Kiểm định Thống kê}:
Kiểm định ghép cặp (Bảng~\ref{tab:statistical_tests}) chỉ ra rằng chênh lệch giữa Section-Aware và Single-Vector là $\bar{\Delta} = +0.0024$ trên nDCG@10 với khoảng tin cậy Bootstrap 95\% CI là $[-0.0296, +0.0333]$ ($t=0.15$, $p$ thô = $0.8811$, $p$ hiệu chỉnh Holm = $1.0000$) và $\bar{\Delta} = +0.0118$ trên MAP với Bootstrap 95\% CI là $[-0.0124, +0.0393]$ ($t=0.88$, $p$ thô = $0.3819$, $p$ hiệu chỉnh Holm = $1.0000$). Vì cả hai khoảng tin cậy đều bao trùm giá trị $0$, chúng tôi kết luận rằng \textbf{Section-Aware Dense đạt sự tương đương thống kê (statistical parity) với Single-Vector Dense trên các vị trí xếp hạng hàng đầu}. Lợi ích thực tế cốt lõi của Section-Aware nằm ở khía cạnh công trình: triệt tiêu hiện tượng cắt xén văn bản ($>95\%$ tránh truncation) và cung cấp các vector phân mục phục vụ tính toán giải thích minh bạch cho nhà tuyển dụng.

\noindent\textbf{3. So sánh với các Phương pháp Chia lát (Chunking)}:
Phương pháp Chunk-Mean đạt nDCG@10 là $0.6193$ ($\bar{\Delta} = +0.0047$, 95\% CI: $[-0.0310, +0.0397]$, $p$ hiệu chỉnh = $1.0000$), trong khi Chunk-MaxSim đạt $0.5992$ ($\bar{\Delta} = -0.0154$, 95\% CI: $[-0.0419, +0.0095]$, $p$ hiệu chỉnh = $1.0000$). Cả hai khoảng tin cậy đều cắt qua $0$, cho thấy việc chia lát cố định không mang lại ưu thế có ý nghĩa thống kê so với nén toàn văn.

\noindent\textbf{4. Tác động trên Mô hình Ngữ cảnh Dài (\textsc{bge-m3})}:
Đối với mô hình BGE-M3 (hỗ trợ $8.192$ tokens), việc phân tách đa vector giúp nâng nDCG@10 từ $0.5937$ lên $0.6086$ ($\bar{\Delta} = +0.0149$, 95\% CI: $[-0.0202, +0.0496]$, $p$ hiệu chỉnh = $1.0000$), không vượt ngưỡng ý nghĩa thống kê chính thức. Mặc dù chỉ số phụ MRR tăng từ $0.6555$ lên $0.7126$ ($+0.0571$), chúng tôi xem đây là một quan sát mang tính khám phá (exploratory) thay vì một khẳng định nhân quả chắc chắn.

\subsection{Trả lời RQ2: Tương quan giữa Từ vựng, Vector Đặc và Hợp nhất Lai}
\label{sec:rq2}
RQ2 phân tích sự phối hợp giữa phương pháp từ vựng (BM25), vector đặc và các cơ chế hợp nhất.
\begin{itemize}
    \item \textbf{Vector Đặc so với BM25}: Vector đặc vượt trội BM25 trên nDCG@10 ($0.6146$ so với $0.5773$ trên Setting A; $0.7153$ so với $0.6720$ trên Setting B). Kiểm định ghép cặp cho thấy giá trị $\bar{\Delta} = +0.0373$ với khoảng tin cậy dương $[+0.0018, +0.0736]$ và $p$ thô = $0.0480$. Tuy nhiên, sau khi hiệu chỉnh đa kiểm định Holm, giá trị $p$ hiệu chỉnh tăng lên $0.5762$. Do đó, dù xu hướng thực nghiệm ưu ái vector đặc nhờ khả năng vượt qua bất đối xứng từ vựng, sự vượt trội chưa đạt mức ý nghĩa thống kê khắt khe dưới họ 15 phép thử.
    \item \textbf{Hợp nhất Dense + BM25}: Mô hình kết hợp chuẩn hóa điểm đạt nDCG@10 là $0.6123$ và MAP là $0.6337$ ($\bar{\Delta} = +0.0127$ so với Single-Vector, 95\% CI: $[-0.0181, +0.0476]$, $p = 0.4566$), đồng thời nâng hệ số tương quan hạng Spearman từ $0.209$ (BM25 thuần) lên $0.276$.
    \item \textbf{Mô hình Không Tham số RRF}: Thuật toán Reciprocal Rank Fusion ($k=60$) kết hợp taxonomy đạt nDCG@10 là $0.6107$ và MAP là $0.6299$ trên Setting A ($0.7109$ và $0.7332$ trên Setting B). Chênh lệch MAP so với mô hình Dense+BM25 được dò tìm tham số là không đáng kể ($\bar{\Delta} = -0.0039$, 95\% CI: $[-0.0233, +0.0119]$, $p = 0.6728$). Nhờ không đòi hỏi điều chỉnh siêu tham số, RRF là lựa chọn thay thế cực kỳ an toàn trong môi trường sản xuất thực tế nhằm phòng chống rủi ro quá khớp.
\end{itemize}

\subsection{Trả lời RQ3: Vai trò của Cây Phân loại Kỹ năng 3 Tầng}
\label{sec:rq3}
RQ3 đánh giá tác động định lượng của các luật ontology kỹ năng. Bảng~\ref{tab:taxonomy_ablation} trình bày ma trận phân rã hai chiều giữa hàm tính điểm (Jaccard Đối xứng so với Độ phủ Bất đối xứng) và các tầng quan hệ ($L_1, L_1+L_2, L_1+L_2+L_3$).

% Bảng phân rã ablation 2 chiều của Taxonomy kỹ năng
\begin{table*}[t]
\centering
\caption{Ma trận phân rã Taxonomy hai chiều: Hàm tính điểm $\times$ Các tầng quan hệ, kiểm định thống kê Holm-Bonferroni so với mốc cơ sở $L_0$ (Dense + BM25).}
\label{tab:taxonomy_ablation}
\resizebox{\textwidth}{!}{%
\begin{tabular}{llrrrrrc}
\hline
\textbf{Tầng quan hệ} & \textbf{Hàm tính điểm} & \textbf{nDCG@10} & \textbf{MAP} & \textbf{$\Delta$ MAP} & \textbf{$p$ thô} & \textbf{$p$ hiệu chỉnh} & \textbf{Ý nghĩa ($\alpha=0.05$)} \\
\hline
Cơ sở $L_0$ & Dense + BM25 (Không Taxonomy) & 0.6123 & 0.6337 & -- & -- & -- & -- \\
Tầng $L_1$ & Jaccard Đối xứng & 0.5910 & 0.5977 & -0.0360 & 0.1651 & 0.6393 & Không \\
Tầng $L_1+L_2$ & Jaccard Đối xứng & 0.5913 & 0.5959 & -0.0378 & 0.1542 & 0.6393 & Không \\
Tầng $L_1$ & Độ phủ Bất đối xứng & 0.5773 & 0.5888 & -0.0449 & 0.1077 & 0.6393 & Không \\
Tầng $L_1+L_2$ & Độ phủ Bất đối xứng & 0.5783 & 0.5896 & -0.0441 & 0.1080 & 0.6393 & Không \\
Tầng $L_1+L_2+L_3$ & Độ phủ Bất đối xứng (Must/Nice) & 0.5775 & 0.5892 & -0.0446 & 0.1066 & 0.6393 & Không \\
Calibrated Hybrid & 0.60 Truy xuất + 0.40 Độ phủ & 0.6162 & 0.6366 & +0.0029 & 0.5430 & 0.6393 & Không \\
\hline
\end{tabular}%
}
\end{table*}


\noindent\textbf{1. Taxonomy Đơn lẻ so với Mô hình Lai}:
Khi hoạt động độc lập không có vector ngữ nghĩa, các luật taxonomy chỉ đạt MAP từ $0.5888$ đến $0.5977$, thấp hơn đáng kể so với các mô hình vector đặc. Điều này khẳng định rằng các luật cứng từ điển không thể thay thế việc hiểu ngữ cảnh tài liệu tự nhiên.

\noindent\textbf{2. Tích hợp Lai và Kiểm định Thống kê}:
Khi tích hợp độ phủ taxonomy vào mô hình Calibrated Hybrid, hệ thống đạt điểm MAP cao nhất trên toàn bộ các mô hình ($0.6366$ trên Setting A; $0.7410$ trên Setting B) và hệ số tương quan thứ hạng cao nhất ($\rho_{\text{mac}} = 0.2954$, $\rho_{\text{mic}} = 0.1944$). Tuy nhiên, khi so sánh với mô hình Dense + BM25 không có taxonomy ($MAP = 0.6337$), mức tăng là $\bar{\Delta} = +0.0029$ trên MAP (95\% CI: $[-0.0055, +0.0127]$, $p$ hiệu chỉnh = $1.0000$) và $\bar{\Delta} = +0.0039$ trên nDCG@10 (95\% CI: $[-0.0070, +0.0157]$, $p$ hiệu chỉnh = $1.0000$). Vì các khoảng tin cậy đều cắt qua $0$, mức cải thiện này không có ý nghĩa thống kê trên tập ứng viên quy mô nhỏ ($N=25$). Mục đích thực tế của việc tích hợp taxonomy là thực thi các điều kiện tiên quyết và xuất danh sách kỹ năng đối sánh phục vụ tính minh bạch trong tuyển dụng.

\subsection{Trả lời RQ4: Hiệu năng Tính toán và Đánh đổi Tái xếp hạng}
\label{sec:rq4}
RQ4 nghiên cứu sự đánh đổi giữa chất lượng xếp hạng và chi phí thời gian thực thi đối với mô hình Cross-Encoder chú ý chéo. Bảng~\ref{tab:two_stage_reranking} mô tả biên Pareto về chất lượng và độ trễ, và Bảng~\ref{tab:latency_breakdown} phân tích chi tiết cấu hình độ trễ chuẩn hóa 3 pha.

% Bảng chất lượng và biên Pareto độ trễ của tái xếp hạng 2 chặng
\begin{table}[t]
\centering
\caption{Biên Pareto giữa Chất lượng xếp hạng và Độ trễ của Mô hình Tái xếp hạng 2 Chặng (Calibrated Hybrid $\rightarrow$ Cross-Encoder) cho pool $N=25$ ứng viên.}
\label{tab:two_stage_reranking}
\resizebox{\columnwidth}{!}{%
\begin{tabular}{lrrrrr}
\hline
\textbf{Kiến trúc / Phương pháp} & \textbf{nDCG@5} & \textbf{nDCG@10} & \textbf{MRR} & \textbf{MAP} & \textbf{$T_{\text{rank}}$ ($N=25$)} \\
\hline
Chặng 1: Calibrated Hybrid Đơn lẻ & 0.5998 & 0.6162 & 0.6764 & 0.6366 & 2.69\,ms \\
Chặng 1+2: Hybrid $\rightarrow$ CE (Top-10) & 0.5513 & 0.5766 & 0.6408 & 0.5767 & 426.9\,ms \\
Chặng 1+2: Hybrid $\rightarrow$ CE (Top-20) & 0.5273 & 0.5532 & 0.6191 & 0.5667 & 853.8\,ms \\
Cross-Encoder Toàn phần (Tất cả $N=25$) & 0.5292 & 0.5531 & 0.6420 & 0.5650 & 1178.0\,ms \\
\hline
\end{tabular}%
}
\end{table}

% Phân tách độ trễ CPU 3 pha (Tiếng Việt)
\begin{table}[tbp]
\centering
\caption{Hồ sơ độ trễ CPU 3 pha chuẩn hóa (Intel Core i7, Pool $N=25$) và đánh đổi chất lượng so với bộ xếp hạng lại Cross-Encoder (Setting~A, $N=71$). $T_{\text{online}} = T_{\text{query}} + T_{\text{rank}}$ là tổng độ trễ trực tuyến đầu-cuối.}
\label{tab:latency_breakdown}
\resizebox{\textwidth}{!}{%
\begin{tabular}{lrrrrrrr}
\toprule
\textbf{Phương pháp / Kiến trúc} & \textbf{nDCG@10} & \textbf{MAP} & \textbf{$T_{\text{index}}$ (ms/tài liệu)} & \textbf{$T_{\text{query}}$ (ms)} & \textbf{$T_{\text{rank}}$ ($N=25$)} & \textbf{$T_{\text{online}}$ (ms)} & \textbf{Đặc tính tính toán} \\
\midrule
BM25 / TF-IDF (Từ khóa) & 0.5773 & 0.5930 & 0.64 & 0.72 & 0.16\,ms & 0.88 & Tích vô hướng chỉ mục ngược thưa \\
Single-Vector Dense (MiniLM) & 0.6146 & 0.6210 & 19.25 & 29.78 & $<0.10$\,ms & 29.79 & Nhân ma trận đặc 384 chiều \\
Section-Aware Dense (Đề xuất) & 0.6170 & 0.6328 & 84.85 & 29.78 & 0.12\,ms & 29.90 & Cosine chéo phân mục động \\
BGE-M3 Flat (8.192 tok) & 0.5937 & 0.6167 & 3012.90 & 1646.12 & $<0.10$\,ms & 1646.13 & Nhân ma trận đặc 1024 chiều \\
BGE-M3 Section-Aware & 0.6086 & 0.6169 & 3314.19 & 1646.12 & 0.13\,ms & 1646.25 & Cosine phân mục 1024 chiều \\
\textbf{Calibrated Hybrid (Đề xuất)} & \textbf{0.6196} & \textbf{0.6351} & 85.49 & 30.75 & \textbf{2.69\,ms} & \textbf{33.44} & Dung hợp MinMax + Taxonomy 3 tầng \\
RRF + Taxonomy ($k=60$) & 0.6169 & 0.6307 & 85.49 & 30.75 & 2.83\,ms & 33.57 & RRF ($k=60$) + độ phủ taxonomy \\
2-Stage: Hybrid $\rightarrow$ CE (Top-10) & 0.5770 & 0.5786 & -- & 30.75 & 426.90\,ms & 457.63 & Tầng 1 + cross-attention Top-10 \\
2-Stage: Hybrid $\rightarrow$ CE (Top-20) & 0.5572 & 0.5681 & -- & 30.75 & 868.02\,ms & 898.77 & Tầng 1 + cross-attention Top-20 \\
Full Cross-Encoder (\texttt{ms-marco}) & 0.5531 & 0.5650 & -- & -- & 1178.05\,ms & 1178.05 & Cross-attention toàn bộ $N=25$ cặp \\
\bottomrule
\end{tabular}%
}
\end{table}


\noindent\textbf{1. So sánh với Cross-Encoder Toàn phần}:
Khi đối đầu với mô hình Cross-Encoder zero-shot (\texttt{ms-marco-MiniLM-L-6-v2}), Calibrated Hybrid cho thấy xu hướng vượt trội trước hiệu chỉnh ($\bar{\Delta} = +0.0632$ trên nDCG@10, 95\% CI: $[+0.0218, +0.1049]$, $p$ thô = $0.0041$; và $\bar{\Delta} = +0.0716$ trên MAP, 95\% CI: $[+0.0270, +0.1196]$, $p$ thô = $0.0039$). Dù khoảng tin cậy Bootstrap hoàn toàn nằm ở miền dương, sau khi áp dụng hiệu chỉnh Holm--Bonferroni trên họ 15 phép so sánh, giá trị $p$ hiệu chỉnh đạt $p_{\text{adj}} = 0.0586$, ngấp nghé trên ngưỡng ý nghĩa danh định $\alpha = 0.05$. Do đó, chúng tôi báo cáo rằng Calibrated Hybrid có xu hướng ưu việt rõ rệt nhưng chưa vượt qua ngưỡng ý nghĩa chính thức sau hiệu chỉnh nghiêm ngặt.

\noindent\textbf{2. Tái xếp hạng Hai Chặng (Cascaded Reranking)}:
Việc sử dụng Cross-Encoder để tái xếp hạng Top-10 ứng viên từ Calibrated Hybrid khiến nDCG@10 giảm từ $0.6162$ xuống $0.5766$ ($\bar{\Delta} = +0.0396$, 95\% CI: $[+0.0061, +0.0770]$, $p$ thô = $0.0342$). Sau hiệu chỉnh Holm, giá trị $p$ đạt $0.4440$, cho thấy chưa có bằng chứng thống kê vững chắc để khẳng định Cross-Encoder làm suy giảm thứ hạng trong kịch bản này.

\noindent\textbf{3. Cấu hình Độ trễ 3 Pha và Tối ưu Hóa Hệ thống}:
Như chi tiết trong Bảng~\ref{tab:latency_breakdown}, độ trễ trực tuyến bị chi phối bởi thời gian mã hóa truy vấn ($T_{\text{query}} \approx 29.8$\,ms cho MiniLM so với $1.646$\,ms cho BGE-M3). Trong khi đó, thời gian tính điểm ứng viên ($T_{\text{rank}}$) cho tập $N=25$ chỉ mất $0.12$\,ms đối với Section-Aware Dense và $2.69$\,ms đối với Calibrated Hybrid. Từ góc độ kỹ thuật hệ thống, khi doanh nghiệp không bắt buộc phải thực thi các luật taxonomy cứng, \textbf{Single-Vector Dense đại diện cho phương án triển khai kinh tế nhất về mặt tính toán} ($T_{\text{online}} \approx 29.8$\,ms, $T_{\text{rank}} < 0.10$\,ms trên CPU), đạt độ chính xác tương đương về mặt thống kê mà không cần duy trì hạ tầng trích xuất đa vector.

\subsection{Kiểm chứng Tổng quát hóa trên Tập Huấn luyện Sạch Tách biệt}
Bảng~\ref{tab:clean_train_benchmark} trình bày kết quả đánh giá trên $76$ JD sạch độc lập thuộc tập huấn luyện ($1.829$ cặp) theo kế hoạch đã cam kết trước (pre-registered plan).

% Bảng thực nghiệm kiểm chứng trên tập Train Clean tách biệt (76 JDs, 1.829 cặp)
\begin{table*}[t]
\centering
\caption{Kiểm chứng Tổng quát hóa trên Tập Huấn luyện Sạch Tách biệt (Held-Out Clean Train Set: 76 JDs, 1.829 cặp tài liệu chưa từng được nhìn thấy khi phát triển luật parser hay khảo sát độ nhạy) qua hai kịch bản Setting A (Toàn bộ 76 truy vấn) và Setting B (71 truy vấn có $\text{rel}\ge 1$).}
\label{tab:clean_train_benchmark}
\resizebox{\textwidth}{!}{%
\begin{tabular}{lrrrrrrrr}
\hline
& \multicolumn{4}{c}{\textbf{Setting A (Toàn bộ 76 Truy vấn)}} & \multicolumn{4}{c}{\textbf{Setting B (Đã lọc $N=71$, $\text{rel}\ge 1$)}} \\
\cmidrule(lr){2-5} \cmidrule(lr){6-9}
\textbf{Mô hình / Kiến trúc} & \textbf{nDCG@10} & \textbf{MRR} & \textbf{MAP} & \textbf{P@5} & \textbf{nDCG@10} & \textbf{MRR} & \textbf{MAP} & \textbf{P@5} \\
\hline
Xếp hạng Ngẫu nhiên (Floor)           & 0.5581 & 0.6553 & 0.5359 & 0.3842 & 0.5974 & 0.7014 & 0.5736 & 0.4113 \\
BM25 / TF-IDF (Từ vựng)                & 0.6376 & 0.7225 & 0.6122 & 0.4395 & 0.6826 & 0.7734 & 0.6554 & 0.4704 \\
MiniLM Đơn Vector Phẳng (Flat)         & 0.6540 & 0.7269 & 0.6503 & 0.4500 & 0.7001 & 0.7781 & 0.6961 & 0.4817 \\
MiniLM Chunk-Mean (128w, stride 64)   & 0.6845 & 0.7725 & 0.6750 & 0.4658 & 0.7328 & 0.8269 & 0.7225 & 0.4986 \\
MiniLM Chunk-MaxSim (ColBERT-style)   & 0.6967 & 0.7827 & 0.6817 & 0.4816 & 0.7457 & 0.8378 & 0.7297 & 0.5155 \\
MiniLM Đa Vector Nhận biết Mục (Ours) & 0.6890 & 0.7492 & 0.6892 & 0.4763 & 0.7376 & 0.8020 & 0.7377 & 0.5099 \\
Dense + BM25 ($L_0$: Không Taxonomy)  & 0.7043 & 0.7794 & 0.6924 & 0.4895 & 0.7539 & 0.8343 & 0.7411 & 0.5239 \\
Calibrated Hybrid (Dense+BM25+Tax)     & 0.6914 & 0.7632 & 0.6822 & 0.4553 & 0.7401 & 0.8170 & 0.7303 & 0.4873 \\
RRF + Taxonomy ($k=60$)                & 0.6227 & 0.6942 & 0.6048 & 0.4605 & 0.6666 & 0.7430 & 0.6474 & 0.4930 \\
Xếp hạng Hoàn hảo Oracle (Ceiling)    & 0.9342 & 0.9342 & 0.9342 & 0.6395 & 1.0000 & 1.0000 & 1.0000 & 0.6845 \\
\hline
\end{tabular}%
}
\end{table*}


Kết quả trên tập Clean Train củng cố vững chắc các phát hiện trên tập Test:
\begin{itemize}
    \item Mô hình Section-Aware Dense đạt nDCG@10 là $0.6890$ và MAP là $0.6892$, nhỉnh hơn Single-Vector Dense ($0.6540$ nDCG@10, $0.6503$ MAP). Kiểm định ghép cặp cho thấy hiệu số $\bar{\Delta} = +0.0351$ (95\% CI: $[+0.0024, +0.0658]$, $p$ thô = $0.0379$). Sau hiệu chỉnh Holm, giá trị $p$ hiệu chỉnh là $0.1517 > 0.05$. Điều này tái khẳng định kết luận: \textbf{sự tương đồng thống kê được bảo toàn nhất quán giữa tập Train và tập Test}.
    \item Mô hình Dense + BM25 đạt kết quả cao nhất trên Setting A ($0.7043$ nDCG@10), trong khi mô hình RRF kém hơn rõ rệt so với mô hình lai có trọng số ($\bar{\Delta} = -0.0687$, $p$ hiệu chỉnh = $0.0003 < 0.05$).
\end{itemize}

\subsection{Thẩm định Chéo Độc lập trên Tập Dữ liệu Mật độ Cao (\textsc{med2425})}
Bảng~\ref{tab:med2425_benchmark} tóm tắt kết quả thử nghiệm trên tập dữ liệu \textsc{med2425} ($41$ JD công nghệ, $277$ CV, $8.575$ cặp, mật độ cạnh tranh cao $\approx 209$ CV/JD).

% Kiểm chứng hai chiều trên med2425 (Tiếng Việt)
\begin{table}[tbp]
\centering
\caption{Kiểm chứng hai chiều trên tập mật độ cao \texttt{med2425/resume-job-fit-merged-v1} (Miền CNTT: $8{,}575$ cặp trên $41$ JD và $277$ CV, nhãn chuyên gia Qwen2.5-32B). Chiều~1 đánh giá Sàng lọc ứng viên ($\text{JD} \rightarrow \text{CV}$, $N=41$, $\approx 209$ CV/JD); Chiều~2 đánh giá Gợi ý việc làm ngược ($\text{CV} \rightarrow \text{JD}$, $N=272$, $\approx 31$ JD/CV).}
\label{tab:med2425_benchmark}
\resizebox{\textwidth}{!}{%
\begin{tabular}{lrrrrrrrrr}
\toprule
 & \multicolumn{5}{c}{\textbf{Chiều 1: Sàng lọc ứng viên ($\text{JD} \rightarrow \text{CV}$, $N=41$)}} & \multicolumn{4}{c}{\textbf{Chiều 2: Gợi ý việc làm ($\text{CV} \rightarrow \text{JD}$, $N=272$)}} \\
\cmidrule(lr){2-6} \cmidrule(lr){7-10}
\textbf{Mô hình / Kiến trúc} & \textbf{nDCG@5} & \textbf{nDCG@10} & \textbf{MRR} & \textbf{MAP} & \textbf{P@5} & \textbf{nDCG@10} & \textbf{MRR} & \textbf{MAP} & \textbf{P@5} \\
\midrule
Xếp hạng ngẫu nhiên (Sàn) & 0.2231 & 0.2350 & 0.5748 & 0.4179 & 0.3854 & 0.3304 & 0.5880 & 0.4650 & 0.4088 \\
BM25 / TF-IDF (Từ khóa) & 0.5736 & 0.5427 & \textbf{0.9129} & 0.6052 & 0.7415 & 0.5728 & 0.8162 & 0.6467 & 0.6471 \\
MiniLM Single-Vector (Flat) & 0.5355 & 0.5393 & 0.8189 & 0.6323 & 0.7366 & 0.5540 & \textbf{0.8324} & 0.6285 & 0.6596 \\
MiniLM Section-Aware (Đề xuất) & 0.5207 & 0.5136 & 0.7859 & 0.6132 & 0.7171 & 0.5725$^\star$ & 0.8211 & 0.6400$^\star$ & 0.6743 \\
Dense + BM25 ($L_0$: Không Tax.) & 0.6198 & 0.5879 & 0.8892 & 0.6438 & 0.7659 & 0.5973$^\star$ & 0.8131 & 0.6648$^\star$ & 0.6846 \\
Calibrated Hybrid (Dense+BM25+Tax) & \textbf{0.6286} & \textbf{0.6249}$^{\star\dagger}$ & 0.9020 & \textbf{0.6833}$^\star$ & \textbf{0.7902} & \textbf{0.6401}$^{\star\dagger}$ & 0.8202 & \textbf{0.6853}$^{\star\dagger}$ & \textbf{0.7000} \\
Xếp hạng Oracle (Trần) & 1.0000 & 1.0000 & 1.0000 & 1.0000 & 0.9951 & 1.0000 & 1.0000 & 1.0000 & 0.8419 \\
\bottomrule
\multicolumn{10}{l}{\footnotesize $^\star$Có ý nghĩa thống kê so với Single-Vector ($p < 0.05$; Chiều~2: Section-Aware Mixed $\bar{\Delta}=+0.0185, p=0.0022$; Hybrid $\bar{\Delta}=+0.0860, p<0.0001$). $^\dagger$Có ý nghĩa so với Dense+BM25 ($p < 0.05$).}
\end{tabular}%
}
\end{table}


Trong bối cảnh mật độ ứng viên dày đặc gấp 8 lần, vai trò của động cơ lai và taxonomy trở nên đặc biệt rõ rệt:
\begin{itemize}
    \item Mô hình Calibrated Hybrid đạt nDCG@10 là \textbf{$0.6242$}, tạo bước nhảy vọt \textbf{$+0.0849$} (+15.7\% tương đối) so với Single-Vector Dense ($0.5393$) và $+0.0363$ so với Dense + BM25 ($0.5879$).
    \item Sự kết hợp giữa vector đặc, từ vựng và bộ lọc cây phân loại giúp giải quyết triệt để bài toán phân loại nhân tài khi có hàng trăm hồ sơ cùng nộp vào một vị trí, chứng minh giá trị thực tiễn to lớn của việc tích hợp tri thức nghiệp vụ trong các tập dữ liệu thực tế quy mô lớn.
\end{itemize}

\subsection{Kiểm toán Chống Rò rỉ Trùng lặp Văn bản}
Để kiểm tra xem các phát hiện thực nghiệm có bị ảnh hưởng bởi hiện tượng lặp lại tài liệu trong tập dữ liệu gốc hay không, Bảng~\ref{tab:leakage_overlap} so sánh kết quả giữa Tập Test Đầy đủ ($477$ CV, bao gồm $68$ CV có xuất hiện trong tập huấn luyện nhưng ghép với các JD khác) và Tập Test Sạch Tách biệt ($409$ CV hoàn toàn độc lập).

% Bảng kiểm định rò rỉ dữ liệu trùng lặp giữa Train và Test
\begin{table}[t]
\centering
\caption{Đánh giá Ảnh hưởng Rò rỉ Dữ liệu Trùng lặp: So sánh hiệu năng trên Tập Test Đầy đủ (477 CV, bao gồm 68 CV trùng với tập huấn luyện) so với Tập Test Sạch Tách biệt (409 CV hoàn toàn không trùng lặp).}
\label{tab:leakage_overlap}
\resizebox{\columnwidth}{!}{%
\begin{tabular}{lrrrr}
\hline
\textbf{Phương pháp} & \textbf{nDCG@10 Đầy đủ} & \textbf{nDCG@10 Sạch} & \textbf{MAP Đầy đủ} & \textbf{MAP Sạch} \\
\hline
BM25 / TF-IDF & 0.5773 & 0.5808 & 0.5930 & 0.5999 \\
Đơn Vector Phẳng (MiniLM) & 0.6146 & 0.6199 & 0.6210 & 0.6290 \\
Đa Vector Nhận biết Mục (Ours) & 0.6170 & 0.6354 & 0.6328 & 0.6458 \\
Dense + BM25 (Không Taxonomy) & 0.6123 & 0.6250 & 0.6337 & 0.6443 \\
Mô hình Lai Hiệu chuẩn (Ours) & 0.6162 & 0.6260 & 0.6366 & 0.6461 \\
Cross-Encoder Toàn phần & 0.5531 & 0.5756 & 0.5650 & 0.5821 \\
\hline
\end{tabular}%
}
\end{table}


Kết quả cho thấy thứ tự xếp hạng tương đối giữa các mô hình được bảo toàn nguyên vẹn. Trên tập sạch, Section-Aware thậm chí còn nới rộng khoảng cách dẫn trước Single-Vector ($0.6354$ so với $0.6199$ nDCG@10; và $0.6458$ so với $0.6290$ MAP), khẳng định các kết luận thực nghiệm của nghiên cứu hoàn toàn trung thực và không phải là hệ quả giả tạo của hiện tượng rò rỉ dữ liệu.
